\section{Benchmark Dataset Design}

\subsection{Dataset Objective}

The objective of the benchmark dataset is to evaluate how different prompt
engineering strategies influence Large Language Model performance on algorithmic
problem-solving tasks.

The dataset is designed to represent common algorithm categories encountered
in computer science education and software engineering technical interviews.


\subsection{Problem Selection}

The benchmark consists of LeetCode-style algorithm problems selected from
multiple categories.

Problems are selected based on the following criteria:

\begin{itemize}

\item Representative algorithmic patterns.

\item Clear problem definitions.

\item Available test cases.

\item Well-established optimal solutions.

\item Known time and space complexity.

\end{itemize}


\subsection{Algorithm Categories}

The benchmark includes problems from the following categories:

\begin{itemize}

\item Array and String Processing

\item Hash Table

\item Two Pointers

\item Sliding Window

\item Linked List

\item Stack and Queue

\item Binary Search

\item Tree

\item Graph

\item Dynamic Programming

\item Greedy Algorithms

\item Backtracking

\end{itemize}


\subsection{Difficulty Distribution}

To evaluate performance across different complexity levels, problems are
divided into three difficulty categories:

\begin{table}[h]

\centering

\begin{tabular}{lll}

\toprule

Difficulty & Description & Purpose \\

\midrule

Easy & Basic algorithm patterns & Evaluate fundamental reasoning \\

Medium & Multiple-step reasoning & Evaluate algorithm selection \\

Hard & Complex optimization problems & Evaluate advanced reasoning \\

\bottomrule

\end{tabular}

\caption{Benchmark difficulty distribution}

\end{table}



\subsection{Dataset Size}

The initial benchmark is designed with approximately:

\begin{itemize}

\item 50 problems for initial evaluation.

\item 100+ problems for extended experiments.

\end{itemize}


The benchmark size balances experimental cost and statistical reliability.


\subsection{Problem Metadata}

Each benchmark problem contains:

\begin{itemize}

\item Problem identifier.

\item Problem title.

\item Algorithm category.

\item Difficulty level.

\item Problem description.

\item Constraints.

\item Expected optimal algorithm.

\item Optimal time complexity.

\item Optimal space complexity.

\item Test cases.

\end{itemize}


The dataset structure enables automated evaluation and reproducible
experiments.
